\documentclass[uplatex]{jsreport}
\usepackage{docmute}
\usepackage{myStyle}
\label{chapter19}

\begin{document}
心理学の実践で最もよく用いられるのが，平均値差の検定です。実験計画のところで学んだように，RCTと組み合わせて使うことで，結果としてみられる群平均の差が効果の大きさと考えられるからです。
本講ではまず，二群の平均値差の検定を考えます。対応がある二群の場合は実はひとつの「差の分布」についての判断になるため，これまでと同様の議論ができるのです。
対応がない二群の場合は，線形モデルの最も単純な場合であり，前期で学んだ様々な要因計画は，分散分析という別名で広く知られている。このように，線形モデルに確率的判断の手続きを加えたものが，平均値差の検定であると言える。実践的側面で注意するべき点は，対抗するモデルである帰無仮説がどのようにおかれているかと，モデル比較の結果が出た後それをどのように解釈するか，である。
\chapter{平均値差の検定}
\section{一要因Withinデザイン，2水準}
効果量

シミュレーションして考える
\section{一要因Betweenデザイン，2水準}
\section{要因計画と検定}
線形モデルと検定の話をひっつける
\section{課題}

\bibliographystyle{jecon_jpa}
\bibliography{sbib}
\end{document}
